\section{Android Networking and \texttt{VpnService}}
Android exposes packet interception on non-rooted devices through \texttt{VpnService}. The service
creates a user-space TUN interface and routes device traffic through an application-controlled
forwarding path \citep{android_vpnservice}. This design enables endpoint telemetry without kernel
modification, but it also imposes practical constraints:
\begin{itemize}[leftmargin=*]
  \item a foreground service is required for long-running capture reliability
  \citep{android_foreground_services},
  \item upstream sockets must be protected from tunnel recursion by calling
  \texttt{protect(...)},
  \item packet processing must remain resource-aware to avoid user-visible overhead.
\end{itemize}
Recent Android releases also require foreground-service type declarations and, for apps targeting
Android 14 or newer, corresponding Play Console declarations. VPN-style monitoring must therefore
be treated as both a networking design and a platform-policy integration problem
\citep{android_foreground_service_types}.
These constraints favor lightweight metadata aggregation over deep packet inspection. MANTA
therefore collects packet and flow-level metadata, aggregates it on-device, and avoids payload
storage.

\section{Flow-Based Monitoring and Metadata-Only Detection}
Flow telemetry abstracts packets into structured records keyed by protocol/endpoint tuples and
counter/timing statistics. NetFlow v9 and IPFIX formalize this approach and define reusable fields
for interoperable flow exchange \citep{rfc3954,rfc7011,rfc7012,iana_ipfix}. Compared with payload
inspection, flow representations reduce content exposure, compress observations for efficient
storage and export, and retain behaviorally meaningful signals such as volume asymmetry,
burstiness, novelty, and frequency change.

Recent anomaly-detection systems such as Kitsune and N-BaIoT show that compact flow-like
representations can support intrusion detection when combined with lightweight autoencoder or
ensemble methods \citep{kitsune_2018,n_baiot_2018}. MANTA adopts this metadata-only approach for
encrypted mobile and endpoint traffic. Its additional concern is deployment privacy: the same
features that help detection may also reveal sensitive application or behavior information.

\section{Encrypted Traffic Analysis and Observer Leakage}
Encrypted transport does not eliminate traffic analysis. Website-fingerprinting work has repeatedly
shown that timing, size, and direction patterns can reveal user activity even when payloads are fully
encrypted. Deep Fingerprinting demonstrated that deep-learning attackers can defeat many
lightweight defenses \citep{deep_fingerprinting_2018}. Later systems work shows that efficient
padding schemes remain difficult to deploy robustly in realistic settings \citep{wf_sok_2022}.
Similar risks appear in mobile traffic: partial encrypted traces can support user-activity detection
\citep{cnn_user_activity_2022}, and concurrent-flow relationships can classify mobile applications
from encrypted traffic \citep{gat_mobile_app_2023}.

These attacks motivate the release/observer split used in MANTA. A coarsened or anonymized export
can reduce what a data consumer learns from the released feature representation. It cannot change
what an on-path observer already sees in the original encrypted traffic sequence. This thesis
evaluates those two leakage paths separately and treats attacker success as an empirical measurement
rather than assuming that removing payloads or direct identifiers is sufficient.

\section{Anomaly Detection for Network Security}
Anomaly detection commonly operates in limited-label settings. Unsupervised or weakly supervised
baselines are attractive because they can be trained on traffic that is mostly normal
\citep{chandola_2009}. Isolation Forest remains a practical reference method for tabular anomaly
scoring because it offers a useful performance/complexity trade-off \citep{isolation_forest_2008}.
In operational security settings, however, threshold policy is often as important as raw score
quality. Base rates, alert costs, and false-positive budgets strongly affect practical utility
\citep{sommer_paxson_2010}.

On-device deployment adds further constraints. Inference must be lightweight, resilient under
battery pressure, and explainable enough for triage. MANTA therefore combines statistical,
exported local-model, and TFLite-compatible scoring with lightweight feature-contribution summaries
that can be generated during inference. This follows the broader mobile-ML trend toward local
inference through TensorFlow Lite/LiteRT-style runtimes, where model format, latency, and fallback
behavior are deployment constraints rather than only implementation details \citep{litert_overview}.

\section{Privacy-Preserving Traffic Representation}
Traditional anonymization techniques such as prefix-preserving pseudonymization protect explicit
identifiers but may still leak semantics through traffic structure. PD-PAn argues that preserving
utility while defending against semantic attacks requires attention to both prefix relationships and
distributional characteristics \citep{pd_pan_2023}. MANTA follows that broader view. It does not
only hash identifiers; it evaluates what remains inferable after identifiers and strong destination
hints have been coarsened, bucketed, or removed from reduced feature views.

\section{Privacy-Aware Representation Learning and Collaboration}
Self-supervised and reconstruction-based learning for encrypted traffic has become a promising path
for anomaly detection under limited labels \citep{et_ssl_2025}. Federated-learning research has also
explored how intrusion-detection models can be improved without centralizing raw data
\citep{federated_privacy_ids_2023,privacy_fl_ids_2024}. MANTA uses these ideas in a limited and
measured form. A privacy student is trained from fuller supervision onto a reduced view, and a
federated simulation operates on that learned latent space rather than on raw full-feature tables.

\section{Positioning}
MANTA differs from earlier work in three ways. First, it integrates Android capture, backend triage,
privacy views, privacy attack benchmarks, and thesis-ready evidence generation into one evaluation
framework. Second, it evaluates a public-only corpus with grouped source-aware splits rather than
relying on a small controlled trace or a purely random split. Third, it treats privacy as two
separate empirical questions: what the released representation leaks, and what a passive observer of
encrypted traffic can still infer.
